\creaCapitulo{Implementacion}

Comenzaremos a partir de la implementación secuencial de los problemas. \cite{Mantas2024} Como comentamos anteriormente, no nos centraremos en la resolución de los problemas per se, puesto que ésta se nos da ya hecha y corregida, sino que nuestros esfuerzos irán en implementar los métodos de resolución en ambas formas de paralelismo, y buscar la máxima eficiencia.

El primer paso, por tanto, será realizar una implementación secuencial de los métodos partiendo de los algoritmos que previamente hemos comentado, Runge-Kutta, Adams-Bashford, y Adamd-Bashford-Moulton, a saber, algoritmos \ref{alg:RungeKutta4}, \ref{alg:AdamsBashford}, y \ref{alg:AdamsBashfordMoulton}, respectivamente.

\creaSeccion{Implementación en OpenMP}
\creaSeccion{Implementación en CUDA}
\creaSeccion{Mediciones}